\begin{pastexam}[label={px:ntut-111-prog-01}]{北科大資工所 111 學年度 第 1 題（節錄 1-1-I 小題）}
（程式追蹤；3 分）

\textbf{Problem 1 [25\%, each 2\%, except 1-1-I 3\%]}

(1) Write the output of the following Python code (1-1-I)$\sim$(1-1-III).

（原題的程式有 26 行，定義了 \texttt{B}、\texttt{CF}、\texttt{A}、\texttt{M}、\texttt{K} 五個函式，
1-1-II、1-1-III 和第 (2)(3) 部分考的是 Python 語法。 這裡只收 1-1-I，程式只列出和它有關的部分，
其餘以 \texttt{...} 表示。）
\begin{minted}[fontsize=\small]{python}
...
def K(perm):
    if len(perm)<=1:return [perm]
    r = []
    for i in range(len(perm)):
        others = K(perm[:i]+perm[i+1:])
        for s in others:
            r = r + [perm[i] + s]
    return r

print(K('XYZ'))                      #(Problem 1-1-I)
...
\end{minted}
\end{pastexam}

\begin{pastsoln}
答案：印出 \texttt{['XYZ', 'XZY', 'YXZ', 'YZX', 'ZXY', 'ZYX']}。

\texttt{K(perm)} 回傳字串 \texttt{perm} 的所有排列：長度 $\le 1$ 時只有它自己一種；否則輪流讓第
\texttt{i} 個字元當開頭（\texttt{perm[i]}），把其餘字元（\texttt{perm[:i]+perm[i+1:]}，也就是拿掉第
\texttt{i} 個）的每一種排列接在它後面。 這正是排列的狀態空間樹（資結書第~20.1~節）：第一層決定
第一個字元，第二層決定第二個，走到葉節點就得到一個排列。 列出的順序是這棵樹的葉節點由左到右
的順序。

追蹤 \texttt{K('XYZ')}：
\begin{center}
\small
\begin{tabular}{clll}
\toprule
\texttt{i} & \texttt{perm[i]} & 其餘字元 & \texttt{K(其餘)} 與接上後的結果 \\
\midrule
0 & \texttt{X} & \texttt{'YZ'} & \texttt{['YZ', 'ZY']} $\Rightarrow$ \texttt{'XYZ', 'XZY'} \\
1 & \texttt{Y} & \texttt{'XZ'} & \texttt{['XZ', 'ZX']} $\Rightarrow$ \texttt{'YXZ', 'YZX'} \\
2 & \texttt{Z} & \texttt{'XY'} & \texttt{['XY', 'YX']} $\Rightarrow$ \texttt{'ZXY', 'ZYX'} \\
\bottomrule
\end{tabular}
\end{center}
其中 \texttt{K('YZ')}：\texttt{i = 0} 時 \texttt{'Y'} 接上 \texttt{K('Z') = ['Z']} 得 \texttt{'YZ'}；
\texttt{i = 1} 時 \texttt{'Z'} 接上 \texttt{K('Y') = ['Y']} 得 \texttt{'ZY'}。 另外兩個同理。
\texttt{r = r + [...]} 每次把新的字串接在串列最後，所以順序就是上表由上到下。

\texttt{print} 印出 Python 串列的寫法：方括號、元素之間用逗號加一個空白、字串用單引號，所以輸出是
\begin{center}
\texttt{['XYZ', 'XZY', 'YXZ', 'YZX', 'ZXY', 'ZYX']}
\end{center}
長度 $n$ 的字串有 $n!$ 種排列，這裡 $3! = 6$ 種，而且因為輸入的字元已經依字母排好，輸出也剛好依
字典序排列。
\end{pastsoln}
